\section{Introduction}

Let \(X\) be a closed smooth four-manifold and let \(J\) be a smooth
almost complex structure on \(X\). A smooth real two-form \(\omega\)
is symplectic if it is nondegenerate and closed, meaning
\(d\omega=0\). A real two-form \(\omega\)
\emph{tames} \(J\) if \(\omega(v,Jv)>0\) for every nonzero tangent
vector \(v\). It is \emph{compatible} with \(J\) if, in addition,
\(\omega(Ju,Jv)=\omega(u,v)\). A taming form is automatically
nondegenerate; thus a closed taming form is symplectic.

Donaldson asked whether the existence of a taming symplectic form
on a closed four-manifold implies the existence of a compatible
one \cite[Section~5.2, Question~2]{Donaldson2006}.
This existence question is distinct from his a priori estimate
conjecture for a prescribed-volume equation.
The elementary projection
\(\omega\mapsto \frac12(\omega+\omega(J\,\cdot,J\,\cdot))\)
shows the obstacle: it preserves positivity on complex lines
but generally loses closedness. Our main result answers the
existence question affirmatively.

\begin{theorem}\label{thm:main}
Suppose that a smooth almost complex structure \(J\) on a closed
connected smooth real four-manifold \(X\) is tamed by a symplectic
form. Then there is a smooth symplectic form compatible with \(J\).
\end{theorem}

The almost complex structure in Theorem~\ref{thm:main} is the given
one. The conclusion imposes no condition on the cohomology class of
the compatible form. In particular, it does not assert that an
arbitrary taming class contains a compatible representative.

There is also an intrinsic description of the compatible cone in terms
of the taming cone and the intersection pairing.

\begin{theorem}[Taming and compatible cones]\label{thm:taming-cones}
Let \(X,J\) satisfy the hypotheses of Theorem~\ref{thm:main}, and give
\(X\) the orientation induced by \(J\). Let
\(\mathcal K_J^t,\mathcal K_J^c\subset H^2(X;\R)\) be the classes
represented by taming and compatible symplectic forms, respectively.
Let \(H_J^-\) be the subspace represented by smooth closed real
anti-invariant two-forms. Set
\[
 \mathcal V=(H_J^-)^\perp,\qquad
 \mathcal C=\mathcal K_J^t\cap\mathcal V,
\]
where orthogonality is for the intersection pairing. Then
\begin{align}
 \mathcal K_J^t&=\mathcal C+H_J^-,\label{eq:tame-decomposition}\\
 \mathcal K_J^c&=
 \{a\in\mathcal C:\ a\cdot y>0
   \text{ for every }0\ne y\in\overline{\mathcal C}\}.
 \label{eq:compatible-duality}
\end{align}
The closure is taken in \(\mathcal V\), and the sum in
\eqref{eq:tame-decomposition} is a Minkowski sum.
\end{theorem}

\begin{corollary}\label{cor:taming-cones}
Under the same hypotheses, write \(h_J^-=\dim H_J^-\) and let
\(b_2^+(X)\) be the positive index of the intersection form.
If \(h_J^-=b_2^+(X)-1\), then
\[
 \mathcal K_J^t=\mathcal K_J^c+H_J^-.
\]
In particular, if \(b_2^+(X)=1\), then
\(\mathcal K_J^t=\mathcal K_J^c\): every taming class has a smooth
compatible symplectic representative.
\end{corollary}

The proofs are given in Section~\ref{sec:cones}. They use separation
in a prescribed cohomology class and the current estimates below.
The inclusion
\(\mathcal K_J^c+H_J^-\subseteq\mathcal K_J^t\)
can be strict when \(h_J^-<b_2^+-1\), even if \(H_J^-=0\);
Example~\ref{ex:strict-cones} gives an explicit four-torus.

The integrable case belongs to the theory of compact complex surfaces.
Buchdahl \cite[Theorem~11]{Buchdahl1999} and Lamari
\cite[Corollary~5.7]{Lamari1999} gave classification-free proofs
that even first Betti number implies the existence of a K\"ahler metric.
Lamari's positive exact current on a surface with odd first Betti number
also obstructs a taming form \cite[Theorem~6.1]{Lamari1999}.
Li and Zhang proved the equivalence between the existence of taming
and compatible symplectic forms on complex surfaces
\cite[Theorem~1.2]{LiZhang2009}.

The pseudoholomorphic-curve approach has a different starting point.
For structures on \(\mathbb{CP}^2\) tamed by the standard symplectic
form, Gromov constructed a compatible form by averaging currents
of pseudoholomorphic lines \cite[Section~2.4.A$'$]{Gromov1985}.
Taubes developed integration over curve moduli spaces to obtain
compatibility for a residual subset of the smooth almost complex
structures tamed by a fixed symplectic form when \(b_2^+=1\)
\cite{Taubes2011,Taubes2017}.
Theorem~1 of \cite{Taubes2011} also states preservation of the taming
form's cohomology class when that class is rational.
Li and Zhang extended this approach to every tamed structure on
\(S^2\times S^2\) and
\(\mathbb{CP}^2\#\overline{\mathbb{CP}}{}^2\), and to further
rational manifolds under hypotheses on embedded pseudoholomorphic
spheres \cite[Theorems~1.2 and~1.3]{LiZhang2015}.

An invariant two-form satisfies
\(\alpha(Ju,Jv)=\alpha(u,v)\); an anti-invariant one satisfies
\(\alpha(Ju,Jv)=-\alpha(u,v)\).
Write \(h_J^-=\dim H_J^-\), where \(H_J^-\) is the anti-invariant
cohomology space defined in Theorem~\ref{thm:taming-cones}.
The condition \(h_J^-=b_2^+-1\) occurs in Buchdahl-type approaches
\cite{TWZZ2022,LinZhou2025}, in the compatibility criteria of
Wang, Wang and Zhu \cite[Theorems~4.3 and~5.1]{WangWangZhu2023},
and in the taming corollary of Wang, Zhang, Zheng and Zhu
\cite[Corollary~1.3]{WangZhangZhengZhu2025}.
The latter paper's generalized Monge--Amp\`ere theorem starts with
an almost K\"ahler structure.
Li and Ning's recent study of spaces of K\"ahler and holomorphically
tamed forms allows the integrable complex structure to vary
\cite{LiNing2026}.

The current approach has a classical cone-duality background.
Sullivan related closed forms positive on a field of tangent directions
to the corresponding structure currents \cite{Sullivan1976}.
Harvey and Lawson characterized K\"ahler manifolds through positive
currents \cite{HarveyLawson1983}; Li and Zhang developed the
almost-complex formulation used to compare taming and compatibility
\cite[Section~3]{LiZhang2009}.
In the argument below, separation gives a positive functional vanishing
on closed invariant forms. An explicit elliptic lift and quantitative
current estimates connect this obstruction to the existence
of a taming form.

\subsection*{The argument and its main estimate}

Fix an auxiliary \(J\)-Hermitian metric. If no compatible form
exists, separation produces a nonzero positive current \(P\)
that annihilates closed invariant forms. An elliptic correction
gives \(T=P+Q\), which annihilates every closed two-form, with
\(Q\) initially an anti-invariant distribution. Represent \(P\)
as a positive weighted average of unit complex-line bivectors,
and let \(\mu\) be the resulting measure on \(X\), its trace measure.
Corrected radial test functions give \(\mu(B_r(x))\le Cr^2\).
Regularization by a finite power of the Hodge resolvent then
proves \(Q\in L^2\) and an angular estimate: over pairs of
complex lines in this representation whose base points are
less than \(r\) apart, the integral of their squared difference,
after parallel transport, is \(O(r^4)\).

The first regularized pairing only bounds \(\|Q\|_2\), because
the negative error in this pairing is only uniformly bounded.
To remove that error, we distinguish the points where \(P\)
has positive two-dimensional density. Write
\[
\theta(x)=\lim_{r\downarrow0}r^{-2}\mu(B_r(x)),
\qquad E=\{x:\theta(x)>0\}.
\]
Existence of the limit is part of the argument. The principal
step, Theorem~\ref{thm:splitting}, proves that the restriction
\(P_E=\ind_E P\) is closed. It applies to any closed current
\(P+Q\) with an anti-invariant \(L^2\) correction and the stated
mass and angular bounds; annihilation of all closed forms is
not needed for this step.

The reason the angular estimate suffices is a balance between
two errors. Planar tangent measures show that, over the same
nearby pairs, the squared normal component of the displacement
between base points, relative to the first complex line and divided
by \(r^2\), has integral \(o(r^2)\).
Closedness of \(P+Q\) trades tangential derivatives of a smooth
density cutoff for a product of angular and transverse errors.
The derivative and normalization of the cutoff contribute
\(r^{-3}\). Cauchy--Schwarz therefore bounds the decisive
boundary term by
\[
r^{-3}\,O(r^2)\,o(r)=o(1),
\]
where the two factors are the square roots of the angular and
transverse moments. The \(L^2\) bound controls the correction's
cutoff error separately. Together these estimates prove closedness
of the density restriction without requiring a quantitative rate
for tangent measure convergence.

Once \(P_E\) is closed, the residual
\(T-P_E=\ind_{X\setminus E}P+Q\) is closed as well. Its positive
summand is concentrated on zero-density centers, where the
negative error in the pairing does tend to zero. Pairing this residual
with \(T\), and using the vanishing of the mixed terms, forces
\(Q=0\). The resulting nonzero positive current \(P=T\)
cannot annihilate a closed taming form.

The relative projection estimate in Lemma~\ref{lem:mixing}, for
fixed smooth orthogonal bundle projections, and the closedness
statement in Theorem~\ref{thm:splitting} may be useful separately.
The regularized pairing estimates use the four-dimensional fact
that anti-invariant two-forms are self-dual.
Sections~\ref{sec:currents} and \ref{sec:mass} construct the
separating current and prove its mass bound.
Section~\ref{sec:energy} establishes the regularization and
angular estimates. Section~\ref{sec:density} proves density
existence and the planar tangent description;
Section~\ref{sec:splitting} uses them to control cutoff boundaries
and prove closedness of \(P_E\).
Section~\ref{sec:conclusion} completes the existence argument.
Section~\ref{sec:cones} develops the prescribed-class separation
and proves the cone description.

A triple of closed two-forms on an oriented four-manifold is
\emph{hypersymplectic} if every nonzero real linear combination is
symplectic and induces the given orientation. The companion
\cite[Theorem~1.1]{OpenAIHypersymplectic2026} combines the current
estimates above with additional prescribed-class and deformation
arguments to deform hypersymplectic triples normalized by
\(\int_X\omega_i\wedge\omega_j=\delta_{ij}\) on closed connected oriented
smooth four-manifolds to hyperk\"ahler triples, while fixing each class
\([\omega_i]\in H^2(X;\R)\).
